\documentclass{report}
\usepackage[a4paper, top=1cm, bottom=1cm, left=2cm, right=2cm, heightrounded]{geometry}

\usepackage{amsmath}
\usepackage[colorlinks=true, urlcolor=blue, linkcolor=black]{hyperref}

\title{DeltaV Architectural Report}
\author{Gabriel Sippl, Simon Geier}
\date{}

\setlength{\parskip}{0.5em}

\begin{document}
\maketitle
\tableofcontents

\chapter{Introduction}
This document outlines the design choices and requirements for each component and will be updated as the design decisions change. \\
\textit{\textbf{DeltaV}} is a Delta robot designed to perform many different \textit{pick-and-place} operations, focusing on precision and modularity. \\
It will be capable of four degrees of freedom, translation in the $x,y$, and $z$ directions as well as rotation around the vertical axis.
The translation components will be handled by the arms' parallelogram structure, whereas the rotation will be handled by a separate module mounted to the end-effector. \\
For now, we will define some broad goals the robot should be able to achieve. \textit{\textbf{DeltaV}} should be capable of
\begin{enumerate}
  \item picking up, rotating, and placing (paper) letters at predetermined positions,
  \item automatically detecting objects in a certain field and moving them to a storage location, and as a final goal
  \item applying solder paste with a syringe to predefined positions and reflow soldering SMD components.
\end{enumerate}

\subsection{Requirements}


\chapter{Mechanical Design}
\section{Base}


\section{Joints}


\section{Arms}


\section{End-Effector}
The end-effector is the moving plate connected to the end of the arms. It's primary function is to hold the modules which allow the robot to interact with the world around it.

\subsection{Structure}
To allow for high acceleration at lower power cost, the structure's inertia must be kept at a minimum.
Lower inertia allows for higher acceleration at lower forces, since inertia is directly linked to an object's mass.
Thus the main goal for the structural design is minimizing weight, while keeping the end-effector stable enough to hold heavier loads without breaking. \\
A rounded triangular model was chosen, since it's three sides allow seamless integration with the three robot arms, while minimizing material around the edges.


\section{Modules}

\chapter{Electrical Design}
\section{Stepper Motor Driver}
For educational purposes, it was decided to build one of the three required stepper motor drivers from scratch. 
However, due to the high cost of individual components, the other two drivers would be purchased ready-made.

\subsection{Component Selection}
\subsubsection{IRF530NPbF - nMOSFET}
Two H-Bridges with four nMOSFETs each will supply the motor with current. 
The IRF530NPbF HEXFET$^\circledR$ power MOSFET was chosen as it strikes a good balance between cost and functionality. 
It features a drain-to-source breakdown voltage of $100V$. 
This allows a margain of $76V$ for potential voltage spikes induced by the inductive load. 
The continuous drain current of $17A$ and sufficiently low $R_{\text{DS(on)}}=90m\Omega$ ensure that any current spikes will not damage the MOSFETs and heat management stays manageable.
Furthermore, the built-in reverse recovery diode can handle up to $60A$ of pulsed drain current with a diode forward voltage of $1.3V$. 
The gate charge $Q_\text{G}=37nC$ is also low enough to allow fast switching without too many switching losses or delays. 

\subsubsection{UCC27710D - MOSFET Gate Driver}
Instead of using pMOS-Technology for the high-side gate drivers, nMOSFETs will be used due to their lower $R_{\text{DS(on)}}$ and thus better heat management for high-power applications. 
This, however, has the drawback that the gate-to-source voltage $V_{GS}$ needs to be significantly higher than the drain-to-source voltage $V_{DS}$ in order for the nMOSFET to conduct. 
Since $V_{DS}$ will sit at around $0V$, driving the MOSFETs with a $5V$ PWM logic signal directly will result in $V_{GS}\approx 2V$ which doesn't reach the threshold voltage of around $3V$. 
Thus, a dedicated high-side gate driver is required to boost the PWM signal to around $36V$ resulting in $V_{GS}=12V$ which is sufficient to drive the MOSFETs. 
The UCC27710D high-side low-side gate driver was chosen as it is specifically designed for motor drivers and features protection circuitry and automatic deadtime against shorts caused by both high and low-side MOSFETs being on at the same time. 
Furthermore, it is capable of supporting a bootstrap capacitor at the high-side gate to boost the voltage to the required $36V$ above GND. 

\subsubsection{1N5819 - Bootstrap Diode}
To charge the bootstrap capacitor from the UCC27710Ds supply voltage of $12V$ a diode with fast switching capabilities and low forward voltage is needed. 
The 1N5819 Shottky diode has a forward voltage of just $0.9V$, repetitive peak reverse voltage of $40V$ and a repetitive peak forward current of $10A$. 
Thus, it is more than sufficient to charge the bootstrap capacitor at high switching frequencies.

\subsubsection{BAT43 - Gate Diode}
The UCC27710D's datasheet recommends adding a diode from the MOSFET's gate terminal back to the HO/LO pins to allow the gate capacitance to discharge when the MOSFET is switched off. 
The BAT43 Shottky diode supports fast switching and low forward voltage with enough continuous forward current to safely discharge the MOSFET.

\chapter{Software Design}

\end{document}